\section{Introduction}
\label{sec:intro}

Classical logic underlies Kolmogorov's probability.  The axioms
of this probability
specify an arithmetic of chance (non-negativity, normalization,
additivity). Beneath that arithmetic sits a logic: the events
form a Boolean $\sigma$-algebra, closed under complement and
satisfying excluded middle.  The complement rule,
$P(A^{\mathsf{c}}) = 1 - P(A)$, is where the two layers touch.  It is an
arithmetical equation about a logical operation, obeyed by
every assignment of degrees of certainty. In the limiting case 
of all-or-nothing probabilities, it becomes
equivalent to the law of excluded middle.

There is no intrinsic reason why probability should be compatible
with exactly one logic as different situations are described by
different logics.  Intuitionistic logic describes computation,
verification, and finite-precision observation, where a proposition
may be confirmable but not refutable.  Linear logic describes
resources, where using an assumption consumes
it~\cite{girard1987}.  Quantum logic describes measurement, where
propositions fail to distribute.  If probability is, as the
Cox--Jaynes tradition has it, \emph{extended
logic}~\cite{cox1946,jaynes2003}, the calculus of degrees of
certainty over a space of propositions, then the logic of those
propositions is a \emph{parameter} of the theory, and Kolmogorov's
calculus is the value of that parameter at Boolean logic.

Physics has already
introduced a use case for a logic not based in Boolean algebra.  
Quantum probability lives on orthomodular
lattices of projections rather than Boolean
algebras~\cite{birkhoff-vonneumann1936}. Gleason's
theorem~\cite{gleason1957} does for quantum probability what this
paper's representation theorems do here: it identifies the states
compatible with the event logic.
It seems natural to conclude that the Boolean substrate 
is a modeling choice.
In this work, investigate what happens when we change this
parameter. Specifically, we ask
"what does probability become when the logic of events is
\emph{intuitionistic}?".

Intuitionistic logic is the natural first non-Boolean instance to 
study using the proofs-as-programs paradigm. 
The choice of (in our case, intuitionistic) 
logic sits upstream of statistics the
way a type system sits upstream of programs: the same syntax of
operations, a different set of valid transformations, and some
derivations that typecheck classically can either fail or split in two.
This logic's
algebraic semantics is mature order theory, well supported by
\textsf{mathlib}~\cite{mathlib2020}: frames, that is, complete
Heyting algebras, modelled on the lattices of open sets.  The
content of the
original (classical) probability definitions and results is not lost:
complete Boolean algebras are exactly the frames where excluded
middle holds,
so the classical theory embeds as a special case rather than being
replaced.

The choice of logic fixes the event algebra.  What remains is to
fix the arithmetic on top of it, keeping only the axioms that can
be stated without a complement.
A \emph{valuation} on a frame $\Omega$ is a monotone, normalized,
modular map $v : \Omega \to \ENNReal$.  For $v$, modularity means
inclusion--exclusion,
$v(a) + v(b) = v(a \sqcup b) + v(a \sqcap b)$, and no axiom mentions a
complement.  In a Heyting algebra the pseudo-complement $a^{\mathsf{c}}$
satisfies
$a \sqcap a^{\mathsf{c}} = \bot$ but not, in general,
$a \sqcup a^{\mathsf{c}} = \top$.  The complement rule therefore degrades
to a derived \emph{inequality} $v(a) + v(a^{\mathsf{c}}) \le 1$, with a gap
\[
  \slack(a) \;=\; 1 - v(a) - v(a^{\mathsf{c}})
\]
measuring exactly the valuation of the region where excluded middle
fails for $a$.  The term \emph{slack}, and this logical reading,
are introduced in this paper.  The quantity itself is not new: it
is the width of the classical Dempster--Shafer belief/plausibility
interval~\cite{dempster1967,shafer1976}. The two align 
on the frame's Boolean core.
Dempster--Shafer reads that width as an agent's ignorance.  Here it
needs no agent: in
the concrete model where $\Omega$ is the opens of
a space carrying a probability measure $\mu$, the slack of an open
$U$ is precisely $\mu(\partial U)$, the measure of its topological
boundary. The slack of an open set is a structural quantity, not a
measure of ignorance.  The theory that results is recognizably the
\emph{shape} of Dempster--Shafer belief functions, but derived from
the underlying logic rather than postulated.

The statistical \emph{operations} are logic-uniform.  The
conditional is still $v(a \sqcap b)/v(b)$, and the product rule and
Bayes symmetry hold verbatim.  What varies is the set of
\emph{identities} those operations satisfy, and with it the licensed
computations.  The classical case split over $\{A, A^{\mathsf{c}}\}$, and
with it the usual denominator of Bayes' theorem
$P(B) = P(B \mid A)\,P(A) + P(B \mid A^{\mathsf{c}})\,P(A^{\mathsf{c}})$, degrades to
a lower bound whose deficit is a conditional slack.  Marginalization
is exact only over \emph{genuine} partitions, families that
actually join to $\top$.  The complement trick
$P(A) = 1 - P(A^{\mathsf{c}})$ degrades to the interval
$[\,v(A),\, 1 - v(A^{\mathsf{c}})\,]$ of width $\slack(A)$, so point
statistics becomes interval statistics exactly on the events the
logic leaves undecided, and collapses back to points on complemented
ones.  Updating itself forks: Dempster's rule and Bayesian
conditioning, classically indistinguishable, give different
posteriors from the same evidence unless its slack vanishes.

The choice of logic changes the
query being answered: the value of an event $A$ is $P(\Box A)$,
where $\Box A$ is the interior of $A$, i.e. the probability of its
verifiable part.
In the measure model, slack is boundary mass, so a nonatomic
distribution queried on intervals behaves entirely classically,
since an interval's boundary is two points, which carry no mass.
The differences are observable precisely at events with charged
boundaries and at semidecidable hypotheses.

A theory this close to excluded middle must be developed with
unusual care, because classicality is easy to smuggle in: a single
informal case split over $\{A, A^{\mathsf{c}}\}$ silently decides an
undecidable event.  The risk is documented even classically, as
Cox's original theorem required repair after Halpern's
counterexample~\cite{cox1946,halpern1999,vanhorn2003}, and the
constructive setting multiplies the opportunities.  
The direct transcription of the classical
Cox statement within an intuitionistic logic setting 
is not merely unprovable but
\emph{unsatisfiable}.  The sum rule is
irreducible: a disjunction's value is provably not a function of
the values of its disjuncts, so
modularity must be posited rather than derived.
Every appearance of classical logic in this work, including at the
meta-theory level, is tracked and justified at the point it is
used, as with the computable \texttt{Decidable} instance behind the
\texttt{CoxModel} non-vacuity witness of Section~\ref{sec:cox}.

Full excluded middle is not the only stop between plain
intuitionistic logic and Boolean logic. Asking what changes at
an intermediate stop is a natural companion question to our main
inquiry into intuitionistic probability.  Weak excluded middle,
$a^{\mathsf{cc}} \sqcup a^{\mathsf{c}} = \top$, is one such case. Using
our valuation machinery, we can demonstrate that slack collapses
to the regularity defect alone under exactly this condition, dropping
the independent De~Morgan gap.

\paragraph{Contributions.} Our contributions are as follows:

\begin{enumerate}
\item \textbf{The two hinges}
  (Sections~\ref{sec:hinge} and~\ref{sec:ds}).  The complement rule
  holds for \emph{every} valuation iff excluded middle holds in the
  frame, an equivalence isolating the exact content of Van Horn's
  classical negation axiom R3~\cite{vanhorn2003}.  Dempster and
  Bayesian conditioning agree iff the evidence has zero slack;

\item \textbf{Structure theorems}
  (Sections~\ref{sec:representation} and~\ref{sec:ds}).  On finite
  frames a valuation is a classical mass function on points.  In
  general it splits into atomic and diffuse parts, and on its
  Boolean core of regular elements it is a belief function;

\item \textbf{The Cox program, dissected} (Section~\ref{sec:cox},
  Appendix~\ref{app:aczel}).  The product-rule half is independent
  of the logic, proved by an Acz\'el/H\"older generator built from
  scratch.  The sum-rule half (a theorem in classical Cox) admits
  no such derivation here, as shown by a five-element
  counterexample, and must therefore be included among the
  defining axioms of a valuation;

\item \textbf{Two grounding models} (Section~\ref{sec:bridges}).
  Any classical probability measure, read on the opens, is a valuation whose
  slack is boundary mass.  On the Sierpi\'nski frame, a slack-free
  classical readout of halting is uncomputable;

\item \textbf{The monad dichotomy} (Section~\ref{sec:conclusion}).
  Product valuations are unique on chain products and provably
  non-unique on the powerset frame, where Fubini fails.  The
  countable fragment composes associatively and commutatively;

\item \textbf{Weak excluded middle.}
  (Sections~\ref{sec:valuation} and~\ref{sec:bridges}).  The global weak
  excluded middle condition is exactly extremal disconnectedness, the class of
  Stonean spaces;

\item \textbf{Mechanization}.
  Everything above is mechanized in Lean~4
  using \textsf{mathlib}, and is (to our knowledge) the first
  mechanization of non-additive probability over intuitionistic
  logic.
\end{enumerate}
